\subsection{Justificación}
Como único centro del Instituto Politécnico Nacional (IPN) especializado en integración tecnológica en la frontera norte, el Centro de Innovación e Integración de Tecnologías Avanzadas (CIITA) tiene una responsabilidad institucional única en el desarrollo industrial y tecnológico de Ciudad Juárez. Un colapso operativo del CIITA no solo afectaría sus actividades internas, sino que tendría un impacto significativo en la región, incluyendo:

\subsubsection{Impacto Institucional} % Capital humano, I+D, económico

\textbf{Formación de capital humano}

Más de 350 estudiantes anuales dependen de los programas de capacitación en automatización industrial y tecnologías avanzadas ofrecidos por el CIITA \cite{ipn2023}. Estos programas son fundamentales para la formación de profesionales altamente capacitados que contribuyen al crecimiento económico y social de la región. La interrupción de estas actividades afectaría directamente la preparación de estos estudiantes y, por ende, la disponibilidad de talento especializado en la industria local.

\textbf{Investigación y desarrollo (I+D)}

El CIITA mantiene colaboraciones activas con más de 15 maquiladoras locales en sectores estratégicos como el aeroespacial y médico. Estas colaboraciones dependen de los proyectos de innovación desarrollados en el centro \cite{ipn2023}. Un deterioro en la capacidad operativa del CIITA podría afectar la continuidad de estos proyectos, lo que a su vez impactaría la competitividad de las empresas asociadas y su capacidad para innovar.

\textbf{Impacto económico}

La inversión inicial requerida para modernizar la infraestructura del CIITA está estimada en 2.5 millones de MXN, lo que representa solo el 18\% del presupuesto anual actual \cite{ipn2023}. Sin embargo, esta inversión garantizaría la continuidad de proyectos que generan aproximadamente 60 millones de MXN en ingresos anuales, consolidando su papel como motor de desarrollo regional.

En resumen, la modernización de la infraestructura tecnológica del CIITA no solo es necesaria para garantizar su operación continua, sino también para mantener su impacto en la formación de capital humano, la investigación y desarrollo, y la competitividad económica de la región. La inversión requerida es estratégica y justificada considerando los beneficios a largo plazo para el centro y la comunidad a la que sirve.

\subsubsection{Cumplimiento Normativo} % LFPDPPP, ISO 27001


\textbf{ Ley Federal de Protección de Datos Personales en Posesión de Particulares (LFPDPPP)}

El cumplimiento de la Ley Federal de Protección de Datos Personales en Posesión de Particulares (LFPDPPP) es fundamental para garantizar la seguridad y privacidad de la información manejada por el Centro de Innovación e Integración de Tecnologías Avanzadas (CIITA). Esta ley establece los lineamientos necesarios para proteger los datos personales de los usuarios, clientes y colaboradores, lo que es especialmente crítico en un entorno tecnológico como el del CIITA, donde se maneja información sensible \cite{lfpdppp}.

\textbf{Implicaciones Financieras}

El incumplimiento de la LFPDPPP puede resultar en sanciones económicas significativas. Las multas por violaciones a esta ley pueden alcanzar hasta 320,000 Unidades de Medida y Actualización (UMA), lo que equivale a millones de pesos, dependiendo de la gravedad de la infracción \cite{lfpdppp}. Además, las sanciones no se limitan a multas económicas; también pueden incluir la suspensión de actividades o la revocación de licencias, lo que afectaría directamente la operación del CIITA.

\textbf{Riesgo Reputacional}

Un incumplimiento de la LFPDPPP no solo tiene implicaciones financieras, sino también reputacionales. La pérdida de confianza por parte de los usuarios, clientes y socios estratégicos puede resultar en una disminución de la colaboración y el financiamiento, lo que a su vez afectaría la sostenibilidad del centro. La reputación del CIITA como un centro de innovación y desarrollo tecnológico es fundamental para mantener las colaboraciones existentes y atraer nuevas oportunidades de negocio.

